\documentclass[twocolumn]{searchchev_preprint}
\usepackage{graphicx}
\begin{document}

\title{Variable stars hidden from Gaia DR3: a sealed-validation search of 13.9 million stars in ZTF,\\
with a pre-registered prediction for Gaia DR4}

\correspondingauthor{Deepnil Ray, \texttt{deepnilray\_ug2025@ashoka.edu.in}}
\author[0000-0001-8779-0700]{Deepnil Ray}
\affiliation{Ashoka University, Sonipat, Haryana 131029, India}
\affiliation{The NoRCEL Institute, United Kingdom}
\email{deepnilray\_ug2025@ashoka.edu.in}

\begin{abstract}
Gaia DR3 classified about 10.5 million of its 1.8 billion sources as variable. I
scored Zwicky Transient Facility (ZTF) DR23 light curves of 13.9 million stars that DR3 left without any
variability result, over 10,700 deg$^2$ at Galactic latitude $|b|\ge15^\circ$. The detector was trained and its threshold
fixed on training tiles, and its performance was measured on validation tiles sealed until then (AUC 0.954). Every
candidate had to show the same period in at least two separate ZTF light curves and, where at least two observing
seasons could be tested, in at least half of them. The search yields 5,778 periodic variables. The ZTF SCoPe classifications, built
independently from the same survey, already list 86.8\% of them as variable with a matching class or period, an
external confirmation that the method finds real variables. Thirty-two stars are listed as variable in no published
catalogue and pass sky-image, contamination and light-curve checks; they are new, and include three detached
eclipsing binaries (orbital periods 1.89, 3.02 and 25.4 d) and two contact binaries. The variables Gaia DR3 missed
are mostly K and M dwarfs with amplitudes near 0.1 mag. Before Gaia DR4 (2 December 2026) I fix and publish a
prediction of what DR4 will report for these stars against matched controls, so DR4 becomes an independent test of
the method and a measurement of which variables Gaia's pipeline misses.
\end{abstract}

\keywords{Variable stars (1761) --- Sky surveys (1464) --- Astronomy data analysis (1858) --- Eclipsing binary stars (444)}

\section{Introduction}
Gaia \citep{Gaia2016} DR3 \citep{Gaia2023} contains the largest homogeneous census of variable stars to date, but its
variability processing only reports sources it identified as variable \citep{Eyer2023,Rimoldini2023}. A star absent from the variability tables may be constant,
or it may be a real variable that the pipeline missed because of Gaia's sparse, scanning-law-driven sampling, its
amplitude, or its brightness. How many real variables are missed, and of which kinds, is not well known.
\citet{Everall2022} describe a selection function for Gaia's (DR2) variable samples as ``hugely valuable'' and note that
building one requires knowing which sources Gaia would have classified as variable had they been variable.

A ground-based survey with dense, multi-year sampling can answer this directly: search the stars Gaia called quiet
and count the real variables among them. ZTF \citep{Bellm2019,Masci2019} has observed the northern sky every few
nights since 2018, so its light curves are well suited. Large ZTF variable catalogues already exist
\citep{Chen2020,Ofek2020,Healy2024}, so the question here is different: not how many variables ZTF contains, but
how many real variables Gaia DR3 left unflagged, and which ones.

The search was designed so that its result can be trusted without trusting the author: the detector and its threshold
were frozen and validated on sealed data before the search, every decision rule for the final catalogue was written
down before the catalogue was assembled, and the catalogue and a prediction for Gaia DR4 are deposited publicly
before DR4 appears. This paper describes the method, the census and the prediction.

\section{Data}
\subsection{Targets}
I selected Gaia DR3 sources with $13\le G\le18.5$, no entry in \texttt{vari\_summary},
\texttt{phot\_variable\_flag} not equal to VARIABLE, and RUWE $<1.4$, at $|b|\ge15^\circ$ and $\delta\ge-25^\circ$. The
sky was divided into the HEALPix partitions of the ZTF HATS archive (orders 4, 5 and 6; 13.4, 3.36 and 0.84 deg$^2$).
At the catalogue freeze, 2,796 partitions (136, 2,646 and 14 of the three orders; 10,700 deg$^2$) had been searched.
They contain 40.4 million such stars. Of these, 13.9 million have an $r$-band light curve with at least 60 good epochs
and enough scatter to pass a pre-filter that keeps 95\% of known variables in the sealed validation set, and were
scored by the detector. Vetting (\S3.3) has so far been run in the 2,646 order-5 partitions (8,880 deg$^2$); the 920
detections in the order-4 and order-6 partitions await vetting and are not in the catalogue.

\subsection{Light curves}
ZTF DR23 PSF-photometry light curves were read from the IRSA HATS archive. Each Gaia star was matched to the ZTF
object with the most epochs within $1.5''$, keeping epochs with no quality flag; stars with another target star
within $3''$ were dropped. The detector runs on the star's $r$-band light curve.
Vetting then uses every $g$, $r$ and $i$ light curve of the star (different fields, CCD quadrants and filters), kept
separate so that each can confirm the others independently.

\section{Method}
\subsection{Detector}
For each light curve I compute shape features: Lomb--Scargle \citep{Lomb1976,Scargle1982} power and false-alarm
probability, amplitude over photometric error, reduced $\chi^2$ about the mean, skewness, the von Neumann ratio and a
box-fit eclipse channel for detached eclipsing binaries. A gradient-boosted decision-tree classifier maps these features
to a score.

\subsection{Sealed validation}
Before any light curve was downloaded, a seeded manifest assigned every tile to training, validation or search. The
classifier was trained on Gaia DR3 classified variables (positives) and Gaia-quiet stars (negatives) from training
tiles only (positives restricted to classes ZTF can see: eclipsing, RR Lyrae, Cepheid, long-period, RS CVn, young
stellar and cataclysmic variables), and its threshold was set from out-of-fold scores on training data. Both were then frozen, and only then
were the validation tiles opened, once. On 356 Gaia-classified variables and 4,000 Gaia-quiet stars in the validation
tiles the classifier reached AUC 0.954 (lower 95\% bound 0.939), against 0.494 for a shuffled-label control. The score
threshold together with a passing periodic or eclipse channel fires on 9 of the 4,000 quiet stars (0.23\%) and
recovers 63\% of the sealed Gaia variables; in the search, a detection must in addition be recovered independently in
each half of the light curve. The 9 quiet stars that fire need not be false positives, because some of those
``negatives'' are unrecognised variables of exactly the kind this search is designed to find. A separate channel
searches periods of 20 min to 2.4 h, below the detector's 0.1-d floor; it contributes 38 catalogue stars.

\subsection{Vetting rules}
Each detection was first checked within $5''$ against 17 published variable-star catalogues; the 16,737 that
matched were not vetted further and are not in the catalogue. The others become candidates only if (i) at least two
of the star's separate ZTF light curves (the detection light curve and others from different fields, CCD quadrants or
filters) each give the same period (within 1\%, factors 0.5, 1 or 2) with $\log_{10}{\rm FAP}<-3$, and no Gaia star
brighter than $G+1$ lies between $1''$ and $8''$; (ii) where at least two observing seasons can be tested (periods up to
60 d), the signal is detected in at least half, and at least two, of them;
(iii) neither the period nor twice it lies within 1\% of a sampling alias (1/3, 1/2, 1 and 2 d, half and one lunar
month, half and one year); (iv) for periods above 1 d, the phases are evenly covered; and (v) no brighter Gaia star
lies within $6''$, and none two magnitudes brighter within $12''$. A star failing any of these, or sharing its period
with other stars in the same tile, is graded C and never claimed. The others are graded A (amplitude $\ge0.08$ mag
with at least three agreeing light curves, including both $g$ and $r$, or amplitude $\ge0.15$ mag with two) or B. For
108 catalogue stars, most with periods above 60 d, fewer than two seasons could be tested. Only grades A
and B are used below.

\subsection{Pre-registered decision rules}
The rules that decide what is claimed were committed publicly on 24 November 2026, before the catalogue was frozen
(see the repository history): when the catalogue freezes; that a star is known, and dropped, if it appears
in VSX \citep{Watson2006} within $5''$ (or within $30''$ with the same period), in any of 17 VizieR variable-star
catalogues, in SIMBAD with a variable-star type, in ALeRCE \citep{SanchezSaez2021} with a periodic class at probability
$\ge0.8$, or in the latest SCoPe release \citep{Healy2024} with variable probability $\ge0.8$ and either a class
probability $\ge0.8$ or a matching period (factors 0.5, 1 or 2, within 1\%); and that every surviving star is inspected on
sky images. A check that fails is retried, and a star it still cannot check is not claimed.

\subsection{Contamination checks}
Every star that no catalogue lists as variable was inspected on Pan-STARRS1 \citep{Chambers2016} $gri$ images
($40''$ cutouts). Galaxies and blends within $3''$ were dropped. Two further checks were applied before claiming a star:
a known variable within $30''$, whose light could leak into the ZTF photometry, or ZTF magnitudes more than 0.5 mag
brighter than Gaia $G$ with no Gaia $G_{BP}-G_{RP}$, a sign that the ZTF source is a blend that Gaia resolves. Either
moves a star to ``uncertain''. Finally every light curve was inspected by eye, and only those with a coherent shape
well above the scatter in both bands are claimed.

\subsection{Use of AI tools}
The analysis code was developed, and the analysis and drafting of this manuscript were carried out, with the
assistance of an AI system (Claude, Anthropic). The author designed and directed the work, checked every number
against the data files and code, and takes full responsibility for the content.

\section{Results}
\subsection{The census}
The search produced 5,778 grade A/B periodic variables among stars Gaia DR3 called quiet (Table~\ref{tab:res},
Figure~\ref{fig:sky}). Their periods run from 0.035 to 351 d (median 1.4 d) and their $r$-band amplitudes, measured
as the 5--95\% range of the light curve, have a median of 0.12 mag (Figure~\ref{fig:pamp}).

\begin{deluxetable}{lc}
\tablecaption{Census at the catalogue freeze (27 November 2026, 03:45 UTC).\label{tab:res}}
\tablehead{\colhead{Quantity} & \colhead{Value}}
\startdata
Sky partitions searched ($|b|\ge15^\circ$) & 2,796 (10,700 deg$^2$) \\
\quad vetted so far (order 5) & 2,646 (8,880 deg$^2$) \\
Gaia-DR3-quiet stars, $13\le G\le18.5$ & 40.4 million \\
\quad scored by the detector & 13.9 million \\
Validated A/B periodic variables & 5,778 \\
\quad in SCoPe as variable, matching class or period & 5,000 (86.8\%) \\
\quad in SCoPe below its variable threshold & 676 (11.7\%) \\
\quad absent from SCoPe & 84 (1.5\%) \\
\quad not in SCoPe's sky coverage & 18 \\
Known in any catalogue & 5,053 \\
Listed as variable in no catalogue & 90 \\
\quad dropped (blend) & 2 \\
\quad held back (possible blend or contamination) & 32 \\
\quad light curve too weak to claim & 24 \\
\quad \textbf{new variable stars} & \textbf{32} \\
\enddata
\end{deluxetable}

\begin{figure*}
\centering
\includegraphics[width=0.82\textwidth]{fig_sky.pdf}
\caption{Sky positions (Galactic coordinates, Mollweide projection) of the 5,778 variables Gaia DR3 did not flag.
Grey: already classified as variable by another catalogue; blue: in no catalogue as a confirmed variable; orange stars:
the 32 new variables. The searched tiles lie at $|b|\ge15^\circ$ and $\delta\ge-25^\circ$.\label{fig:sky}}
\end{figure*}

\subsection{An external check: SCoPe}
SCoPe \citep{Healy2024} classifies ZTF light curves with its own features and models, independently of this
work; I used its latest public release (Zenodo record 20057104). Of the 5,760 candidates inside its coverage, it lists 5,000 (86.8\%) as variable with a matching class or
period. A detector that found noise or artefacts would not agree with an independent classifier at this rate, so this is
the strongest single check that the 5,778 detections are real. A further 676 (11.7\%) are in SCoPe but below its own
variable threshold, or with a different period and class; for these the present work supplies a vetted period
confirmed in independent light curves and across seasons. They are reported as independently recovered, not as
discoveries.

\subsection{New variable stars}
Ninety stars are listed as variable in no catalogue. After the image, contamination and light-curve checks, 32 are
claimed as new (Table~\ref{tab:new}, Figure~\ref{fig:lc}). They include three detached eclipsing binaries: Gaia DR3
1198491283713804928 ($P=1.890$ d), 1052374297278903936 ($P=3.024$ d) and 346276760375634816, whose 25.4-d orbit with
narrow eclipses makes it a rare long-period system and a good target for radial-velocity follow-up. Two are contact
binaries with periods of 0.248 and 0.290 d. The others are periodic variables at 0.24 to 22 d with amplitudes of
0.06--0.32 mag; their types are not assigned here, and for these the true period may be twice the ZTF period.
A further 24 stars passed every catalogue and image check but have light curves too weak to claim; they are listed
separately in the Zenodo record as candidates.

\begin{figure*}
\centering
\includegraphics[width=0.95\textwidth]{fig_lc.png}
\caption{Six of the new variables: ZTF $g$ (green, upper panels) and $r$ (red, lower panels) light curves folded at
the ZTF period (left of each pair) and twice it (right). From top left: the 25.4-d eclipsing binary
346276760375634816; the eclipsing binaries 1052374297278903936 and 1198491283713804928; the 22-d variable
321766687685841536; the 0.34-d variable 220985864201130880; and the 0.248-d contact binary 4230919197495489024.\label{fig:lc}}
\end{figure*}

\subsection{Which variables Gaia DR3 missed}
Two properties of the sample describe what Gaia's pipeline missed, and a third check is reported for completeness.

\emph{They are mostly red dwarfs.} Of the 4,450 candidates with parallax over error above 5, 62\% have
$G_{BP}-G_{RP}>1.8$ and the median absolute magnitude is $M_G=8.1$: almost all lie on the main sequence, and most are
K and M dwarfs (Figure~\ref{fig:cmd}). Among these, the stars that no catalogue had classified are redder still (median
$G_{BP}-G_{RP}=2.39$ against 2.05 for the known ones).

\emph{They have small amplitudes.} The median $r$-band amplitude is 0.12 mag, and 0.097 mag for the stars no catalogue had
classified; these are also longer in period (median 2.3 d against 1.3 d). Low-amplitude, longer-period variability
of the kind produced by starspots on cool dwarfs is what both Gaia and earlier ZTF catalogues missed most.

\emph{Gaia observations.} For 391 candidates paired with a control star of matching brightness and colour within
$0.5^\circ$, the number of Gaia DR3 $G$-band observations (\texttt{phot\_g\_n\_obs}) is the same (median paired
difference $-3$, Wilcoxon $p=0.33$; Figure~\ref{fig:nobs}). Because pairs this close share Gaia's scanning pattern,
this test cannot show whether sampling contributed to the omission; that needs a comparison across sky regions with
different scan coverage, which is planned with the DR4 test.

These are descriptive results from the census alone. A completeness-corrected selection function, from injecting
synthetic variables into real ZTF light curves, will be presented with the DR4 test.

\begin{figure}
\centering
\includegraphics[width=\columnwidth]{fig_pamp.pdf}
\caption{Period and amplitude of the 5,778 variables (colours as in Figure~\ref{fig:sky}).\label{fig:pamp}}
\end{figure}

\begin{figure}
\centering
\includegraphics[width=\columnwidth]{fig_cmd.pdf}
\caption{Gaia colour--absolute magnitude diagram of candidates with parallax over error above 5 (colours as in
Figure~\ref{fig:sky}). Almost all are main-sequence stars, most of them K and M dwarfs.\label{fig:cmd}}
\end{figure}

\begin{figure}
\centering
\includegraphics[width=\columnwidth]{fig_nobs.pdf}
\caption{Number of Gaia DR3 $G$-band observations for 391 candidates and their matched quiet control stars. The
distributions are indistinguishable (Wilcoxon $p=0.33$).\label{fig:nobs}}
\end{figure}

\section{Prediction for Gaia DR4}
The following is fixed before DR4 appears; the files, the scoring script and their SHA-256 hashes are in the Zenodo
record. Stars are matched through DR4's \texttt{dr3\_neighbourhood} table, with a $1''$ position match as a reported
fallback.

\textbf{P1.} DR4 flags a larger fraction of the 5,778 candidates as VARIABLE than of 550 control stars, each chosen at
random among Gaia-quiet stars within $0.5^\circ$ of a candidate with $|\Delta G|<0.25$ and
$|\Delta(G_{BP}-G_{RP})|<0.15$, at $z>3.3$. The POSSIBLE and NON\_VARIABLE fractions are reported separately.

\textbf{P2.} DR4 $G$-band light curves show the ZTF period (factor 0.5, 1 or 2, within 1\%) with false-alarm probability
$<10^{-3}$ more often for candidates than for controls tested at their paired candidate's period.

\textbf{P3.} Among candidates, the DR4 VARIABLE fraction rises with amplitude and falls for faint stars and for periods
beyond about 10 d, measured in bins fixed now (8 in $\log P$ over 0.05--500 d, 5 in amplitude, 1-mag bins in $G$ from 12
to 19, Wilson intervals).

No target fraction is set, and the outcome will be reported whatever it is.

\section{Discussion}
The method is general. Any time-domain survey can be searched for variables a reference catalogue missed in the same
way: freeze the detector on sealed tiles, demand independent-light-curve and season-to-season confirmation, check
prior art under rules written in advance, and test the result against the next data release. The three ingredients
that make the result trustworthy, sealed validation, external confirmation by an independent classifier, and an
advance prediction, cost little and could be added to most variability searches.

Limitations should be stated plainly. The ZTF DR1 candidates of \citet{Ofek2020} could not be obtained and were not
checked, and the ZTF Variability Archive (ZVAR) DR1 \citep{Graham2026}, released after the catalogue freeze, has not
yet been checked. The amplitude used here includes photometric noise, so it overestimates the true amplitude of the faintest
stars. The count of stars scored is taken from the search logs. Completeness has not yet been measured, so the
census is a lower limit on the number of variables Gaia DR3 missed in these tiles. The sky coverage is limited to
$|b|\ge15^\circ$ and $\delta\ge-25^\circ$; crowded low-latitude fields, where blending is worse, are not included.

The search continues on further tiles. The DR4 test, the completeness measurement and a selection function in
period, amplitude, magnitude and Gaia scanning-law coverage will follow in a full paper.

\section*{Data availability}
The frozen catalogue, the new-variable list, the DR4 prediction with its control list and scoring script, the image and
light-curve checks, and a code snapshot are archived on Zenodo (\href{https://doi.org/10.5281/zenodo.24579691}{doi:10.5281/zenodo.24579691}, deposited 27 November 2026). Code and
results are at \url{https://github.com/deepnilray/SearchChev}.

\begin{acknowledgments}
I thank Prof. T. R. Seshadri (Ashoka University) for reviewing this project and for his advice on its method.
This work used data from the Zwicky Transient Facility, obtained through the NASA/IPAC Infrared Science Archive, and
data from the ESA mission Gaia, processed by the Gaia Data Processing and Analysis Consortium. It made use of the
ZTF SCoPe classifications, the ALeRCE broker, the AAVSO VSX catalogue, the SIMBAD and VizieR services of CDS, the
Pan-STARRS1 image archive, and the Gaia mirror of the Leibniz Institute for Astrophysics Potsdam. The use of AI tools is described in the Method section.
\end{acknowledgments}

\section*{Conflicts of interest}
The author declares no conflicts of interest.

\begin{thebibliography}{}
\bibitem[Bellm et al.(2019)]{Bellm2019} Bellm, E. C., et al. 2019, PASP, 131, 018002
\bibitem[Chambers et al.(2016)]{Chambers2016} Chambers, K. C., et al. 2016, arXiv:1612.05560
\bibitem[Chen et al.(2020)]{Chen2020} Chen, X., et al. 2020, ApJS, 249, 18
\bibitem[Everall \& Boubert(2022)]{Everall2022} Everall, A., \& Boubert, D. 2022, MNRAS, 509, 6205
\bibitem[Eyer et al.(2023)]{Eyer2023} Eyer, L., et al. 2023, A\&A, 674, A13
\bibitem[Gaia Collaboration et al.(2016)]{Gaia2016} Gaia Collaboration, Prusti, T., et al. 2016, A\&A, 595, A1
\bibitem[Gaia Collaboration et al.(2023)]{Gaia2023} Gaia Collaboration, Vallenari, A., et al. 2023, A\&A, 674, A1
\bibitem[Graham et al.(2026)]{Graham2026} Graham, M. J., et al. 2026, arXiv:2611.16166
\bibitem[Healy et al.(2024)]{Healy2024} Healy, B. F., et al. 2024, ApJS, 272, 14
\bibitem[Lomb(1976)]{Lomb1976} Lomb, N. R. 1976, Ap\&SS, 39, 447
\bibitem[Masci et al.(2019)]{Masci2019} Masci, F. J., et al. 2019, PASP, 131, 018003
\bibitem[Ofek et al.(2020)]{Ofek2020} Ofek, E. O., et al. 2020, MNRAS, 499, 5782
\bibitem[Rimoldini et al.(2023)]{Rimoldini2023} Rimoldini, L., et al. 2023, A\&A, 674, A14
\bibitem[S\'anchez-S\'aez et al.(2021)]{SanchezSaez2021} S\'anchez-S\'aez, P., et al. 2021, AJ, 161, 141
\bibitem[Scargle(1982)]{Scargle1982} Scargle, J. D. 1982, ApJ, 263, 835
\bibitem[Watson et al.(2006)]{Watson2006} Watson, C. L., Henden, A. A., \& Price, A. 2006, SASS, 25, 47
\end{thebibliography}

\appendix
\section{The new variable stars}
\startlongtable
\begin{deluxetable*}{lrrrrrc}
\tablecaption{The 32 new variable stars. $P$ is the orbital period for EA (detached eclipsing) and EW (contact)
binaries; for the others it is the ZTF period, and the true period may be twice it. Amp.\ is the ZTF $r$-band
5--95\% range.\label{tab:new}}
\tablehead{\colhead{Gaia DR3 source\_id} & \colhead{RA (deg)} & \colhead{Dec (deg)} & \colhead{$G$} & \colhead{$P$ (d)} & \colhead{Amp.} & \colhead{Type}}
\startdata
\input{table.tex}
\enddata
\end{deluxetable*}

\end{document}
